% 2 · Objectives
\section{Objectives}

\begin{frame}{Broad objective: a diagnosis, by scale of impact}
  \vfill
  \begin{center}
    \begin{minipage}{.82\textwidth}
      \large Whether, where, and at what scale AI data centres put pressure on planetary
      boundaries, as a basis for \textbf{responses matched to each scale of impact}.
    \end{minipage}
  \end{center}
  \vfill
  \centering
  \bignum{Whether}{Is the pressure real\\and growing?}\hfill
  \bignum{Where}{Which grids, basins\\and communities?}\hfill
  \bignum{Scale}{Local, regional\\or global?}
  \vfill
  \note{%
    \NoteKey{We diagnose whether, where and at what scale AI data centres press on planetary boundaries.}
    \NoteSlide{The aim, then three questions. \emph{Whether}: is the pressure real and growing? \emph{Where}: which grids,
    basins and communities? \emph{Scale}: local, regional or global?}
    \NoteWhy{The scale decides who has to respond: a town council, a river-basin or grid authority, or a global agreement.}
    \NoteSay{25}{Our broad objective is a diagnosis: whether AI data centres put real pressure on planetary boundaries,
    where that pressure lands, and at what scale, so that responses can be matched to each scale.}
  }
\end{frame}

\begin{frame}{Specific objectives}{}
  {\footnotesize\renewcommand{\arraystretch}{1.35}
  \begin{tabularx}{\textwidth}{@{}>{\bfseries\raggedright\arraybackslash}p{2.7cm}
    >{\raggedright\arraybackslash}p{2.5cm} >{\raggedright\arraybackslash}X@{}}
    \toprule
    & Earth sphere & What we measure \\
    \midrule
    O1\enspace Energy and carbon & Atmosphere
      & Regional capacity, electricity use and grid carbon; growth rates across the three
        phases; why the cloud era stayed flat while the AI era does not \\
    O2\enspace Materials, water and waste & Hydrosphere, lithosphere
      & Water-stress exposure of capacity and nearby people; water use and sources; geography
        of critical minerals and e-waste \\
    O3\enspace Economy, politics and society & Anthroposphere
      & How unevenly capacity is spread relative to population; local opposition;
        supply-chain concentration \\
    O4\enspace Synthesis & All
      & Findings placed on the planetary and safe-and-just boundaries, classified as local,
        regional or global \\
    \bottomrule
  \end{tabularx}}
  \src{Sakschewski et al.\ (2025); Rockström et al.\ (2023) \emph{Nature}}
  \note{%
    \NoteKey{Four objectives, one per part of the Earth system, plus a synthesis.}
    \NoteSlide{O1 energy and carbon (atmosphere). O2 materials, water and waste (hydrosphere = water; lithosphere = rock and
    minerals). O3 economy, politics and society (anthroposphere = the human sphere). O4 synthesis: everything placed on the
    planetary boundaries and sorted into local, regional or global.}
    \NoteWhy{Each later results section answers one of these rows.}
    \NoteSay{30}{We split the work into four objectives. O1 looks at electricity and carbon, and why the cloud era stayed
    flat. O2 looks at water stress, water use, minerals and e-waste. O3 looks at how unevenly capacity is spread and at
    local opposition. O4 puts all of it on the planetary boundaries, by scale.}
    \begin{NoteAsk}
      \item \emph{What are ``safe-and-just'' boundaries?} Rockström et al.\ (2023) add a justice test: limits set so that
        people are not significantly harmed, not only so that Earth stays stable.
    \end{NoteAsk}
  }
\end{frame}
